% ============================================================
%  MTH636 – Game Theory | Lecture 13
%  The Indifference Principle and Computing Nash Equilibria
% ============================================================
\documentclass[12pt, a4paper]{article}

% ============================================================
%  MTH636 -- Game Theory
%  Notes Preamble — Color-Coded, TikZ-Enhanced
% ============================================================

% --- Core packages ---
\usepackage[T1]{fontenc}
\usepackage[utf8]{inputenc}
\usepackage{lmodern}
\usepackage[margin=2.5cm, top=2.8cm, bottom=2.8cm]{geometry}
\usepackage{microtype}

% --- Math ---
\usepackage{amsmath, amssymb, amsthm, mathtools}
\usepackage{bm}        % bold math
\usepackage{cancel}    % strike-through in math

% --- Color & Boxes ---
\usepackage[dvipsnames, svgnames, x11names, table]{xcolor}
\usepackage[most]{tcolorbox}
\tcbuselibrary{breakable, skins, theorems}

% --- TikZ & Diagrams ---
\usepackage{tikz}
\usepackage{tikz-qtree}
\usetikzlibrary{arrows.meta, positioning, shapes.geometric, fit, backgrounds,
                decorations.pathreplacing, calc, matrix, trees}
\usepackage{pgfplots}
\pgfplotsset{compat=1.18}

% --- Tables ---
\usepackage{booktabs}
\usepackage{multirow}
\usepackage{array}
\usepackage{tabularx}
\usepackage{makecell}

% --- Typography & Layout ---
\usepackage{enumitem}
\usepackage{parskip}
\usepackage{fancyhdr}
\usepackage{titlesec}
\usepackage{titletoc}
\usepackage{hyperref}
\usepackage{graphicx}
\usepackage{float}
\usepackage{caption}

% --- Color Palette ---
\definecolor{defblue}{HTML}{1A5276}      % Definition header
\definecolor{defbluebg}{HTML}{D6EAF8}    % Definition background
\definecolor{thmgreen}{HTML}{1E8449}     % Theorem header
\definecolor{thmbg}{HTML}{D5F5E3}        % Theorem background
\definecolor{exorange}{HTML}{B7770D}     % Example header
\definecolor{exbg}{HTML}{FEF9E7}         % Example background
\definecolor{warnred}{HTML}{922B21}      % Warning/Key Insight header
\definecolor{warnbg}{HTML}{FADBD8}       % Warning background
\definecolor{noticegray}{HTML}{2C3E50}   % Note header
\definecolor{noticebg}{HTML}{EBF5FB}     % Note background
\definecolor{taketawayyellow}{HTML}{7D6608} % Key Takeaway header
\definecolor{takeawaybg}{HTML}{FDFEFE}   % Takeaway background
\definecolor{accentpurple}{HTML}{6C3483} % Accent / proof
\definecolor{lightgray}{HTML}{F2F3F4}    % shading

% ============================================================
%  TCOLORBOX ENVIRONMENTS
% ============================================================

% -- Definition (Blue) --
\newtcolorbox{defbox}[2][]{
  enhanced, breakable,
  colback=defbluebg, colframe=defblue, coltitle=white,
  fonttitle=\bfseries\sffamily,
  title={Definition: #2},
  attach boxed title to top left={yshift=-2mm, xshift=4mm},
  boxed title style={colback=defblue, rounded corners},
  arc=4pt, boxrule=1.2pt,
  left=6pt, right=6pt, top=8pt, bottom=6pt,
  #1
}

% -- Theorem / Proposition (Green) --
\newtcolorbox{thmbox}[2][]{
  enhanced, breakable,
  colback=thmbg, colframe=thmgreen, coltitle=white,
  fonttitle=\bfseries\sffamily,
  title={Theorem: #2},
  attach boxed title to top left={yshift=-2mm, xshift=4mm},
  boxed title style={colback=thmgreen, rounded corners},
  arc=4pt, boxrule=1.2pt,
  left=6pt, right=6pt, top=8pt, bottom=6pt,
  #1
}

% -- Proposition (Green variant) --
\newtcolorbox{propbox}[2][]{
  enhanced, breakable,
  colback=thmbg, colframe=thmgreen!70!black, coltitle=white,
  fonttitle=\bfseries\sffamily,
  title={Proposition: #2},
  attach boxed title to top left={yshift=-2mm, xshift=4mm},
  boxed title style={colback=thmgreen!70!black, rounded corners},
  arc=4pt, boxrule=1.2pt,
  left=6pt, right=6pt, top=8pt, bottom=6pt,
  #1
}

% -- Example (Orange/Yellow) --
\newtcolorbox{exbox}[2][]{
  enhanced, breakable,
  colback=exbg, colframe=exorange, coltitle=white,
  fonttitle=\bfseries\sffamily,
  title={Example: #2},
  attach boxed title to top left={yshift=-2mm, xshift=4mm},
  boxed title style={colback=exorange, rounded corners},
  arc=4pt, boxrule=1.2pt,
  left=6pt, right=6pt, top=8pt, bottom=6pt,
  #1
}

% -- Key Insight / Warning (Red) --
\newtcolorbox{warnbox}[2][]{
  enhanced, breakable,
  colback=warnbg, colframe=warnred, coltitle=white,
  fonttitle=\bfseries\sffamily,
  title={Key Insight: #2},
  attach boxed title to top left={yshift=-2mm, xshift=4mm},
  boxed title style={colback=warnred, rounded corners},
  arc=4pt, boxrule=1.2pt,
  left=6pt, right=6pt, top=8pt, bottom=6pt,
  #1
}

% -- Note (Gray-Blue) --
\newtcolorbox{notebox}[1][]{
  enhanced, breakable,
  colback=noticebg, colframe=noticegray!60, coltitle=white,
  fonttitle=\bfseries\sffamily,
  title={Note},
  attach boxed title to top left={yshift=-2mm, xshift=4mm},
  boxed title style={colback=noticegray!60, rounded corners},
  arc=4pt, boxrule=0.8pt,
  left=6pt, right=6pt, top=8pt, bottom=6pt,
  #1
}

% -- Key Takeaway --
\newtcolorbox{takebox}[1][]{
  enhanced, breakable,
  colback=takeawaybg, colframe=taketawayyellow, coltitle=white,
  fonttitle=\bfseries\sffamily,
  title={Key Takeaway},
  attach boxed title to top left={yshift=-2mm, xshift=4mm},
  boxed title style={colback=taketawayyellow, rounded corners},
  arc=4pt, boxrule=1.2pt,
  left=6pt, right=6pt, top=8pt, bottom=6pt,
  #1
}

% -- Intuition --
\newtcolorbox{intubox}[1][]{
  enhanced, breakable,
  colback=accentpurple!8, colframe=accentpurple!60, coltitle=white,
  fonttitle=\bfseries\sffamily,
  title={Intuition},
  attach boxed title to top left={yshift=-2mm, xshift=4mm},
  boxed title style={colback=accentpurple!80, rounded corners},
  arc=4pt, boxrule=0.8pt,
  left=6pt, right=6pt, top=8pt, bottom=6pt,
  #1
}

% ============================================================
%  PAYOFF MATRIX HELPER
%  Usage: inside tabular, use \payoff{a}{b} for (a,b)
% ============================================================
\newcommand{\payoff}[2]{$(#1,\;#2)$}
\newcommand{\nash}[2]{\underline{$#1$},\underline{$#2$}}  % underline NE payoffs
\newcommand{\Nset}{\mathbb{N}}
\newcommand{\Rset}{\mathbb{R}}
\newcommand{\Sset}{\mathcal{S}}
\newcommand{\Ui}{u_i}

% ============================================================
%  SECTION FORMATTING
% ============================================================
\titleformat{\section}
  {\large\bfseries\sffamily\color{defblue}}
  {\thesection}{1em}{}[\titlerule]

\titleformat{\subsection}
  {\normalsize\bfseries\sffamily\color{thmgreen!80!black}}
  {\thesubsection}{1em}{}

\titleformat{\subsubsection}
  {\normalsize\itshape\sffamily\color{exorange!80!black}}
  {\thesubsubsection}{1em}{}

% ============================================================
%  HEADER / FOOTER
% ============================================================
\pagestyle{fancy}
\fancyhf{}
\fancyhead[L]{\sffamily\small\color{defblue} MTH636 -- Game Theory}
\fancyhead[R]{\sffamily\small\color{defblue} \leftmark}
\fancyfoot[C]{\sffamily\small\thepage}
\renewcommand{\headrulewidth}{0.4pt}

% ============================================================
%  HYPERREF SETUP
% ============================================================
\hypersetup{
  colorlinks=true,
  linkcolor=defblue,
  urlcolor=thmgreen,
  citecolor=accentpurple,
  bookmarksnumbered=true,
  pdftitle={MTH636 Game Theory Notes},
  pdfauthor={MTH636 Student Notes}
}

% ============================================================
%  MISC HELPERS
% ============================================================
\newcommand{\BR}{\text{BR}}        % Best Response
\newcommand{\NE}{\text{NE}}        % Nash Equilibrium
\newcommand{\SPNE}{\text{SPNE}}    % Subgame Perfect NE
\newcommand{\Dom}{\text{dom}}      % Dominance
\newcommand{\mixs}[1]{\sigma_{#1}} % Mixed strategy
\newcommand{\pures}[1]{s_{#1}}     % Pure strategy


% ============================================================
\begin{document}

% Title
\begin{center}
  {\LARGE\bfseries\sffamily\color{defblue} MTH636 -- Game Theory}\\[6pt]
  {\Large\sffamily Lecture 13: The Indifference Principle and Finding NE}\\[4pt]
\end{center}
\vspace{0.5cm}
\hrule height 1.5pt
\vspace{0.3cm}

% ============================================================
\section{The Indifference Principle}

\begin{thmbox}{The Indifference Principle}
  Suppose $\sigma^*$ is a Nash Equilibrium. For any player $i$, if there are two pure strategies $s_i^1, s_i^2 \in S_i$ that are played with strictly positive probability (i.e., $\sigma_i^*(s_i^1) > 0$ and $\sigma_i^*(s_i^2) > 0$), then:
  \[
    U_i(s_i^1, \sigma_{-i}^*) = U_i(s_i^2, \sigma_{-i}^*)
  \]
\end{thmbox}

\textbf{Proof Idea:}
Suppose for contradiction that $U_i(s_i^1, \sigma_{-i}^*) > U_i(s_i^2, \sigma_{-i}^*)$. 
Player $i$ could define a new mixed strategy $\bar{\sigma}_i$ by taking all the probability mass assigned to $s_i^2$ and shifting it to $s_i^1$:
\[
  \bar{\sigma}_i(s_i) = \begin{cases} 
    \sigma_i^*(s_i) & \text{if } s_i \neq s_i^1, s_i^2 \\ 
    0 & \text{if } s_i = s_i^2 \\ 
    \sigma_i^*(s_i^1) + \sigma_i^*(s_i^2) & \text{if } s_i = s_i^1 
  \end{cases}
\]
By evaluating the expected utility $U_i(\bar{\sigma}_i, \sigma_{-i}^*)$, one can show that Player $i$ obtains a strictly higher expected payoff than $U_i(\sigma_i^*, \sigma_{-i}^*)$, which contradicts the assumption that $\sigma^*$ is a Nash Equilibrium.

\begin{notebox}
  \textbf{Corollary:} For any Nash Equilibrium $\sigma^*$, the expected utility $U_i(\sigma^*)$ is exactly equal to the expected utility of \textbf{any} pure strategy $s_i$ in its support:
  \[
    U_i(\sigma^*) = U_i(s_i, \sigma_{-i}^*) \quad \forall s_i \in S_i \text{ such that } \sigma_i^*(s_i) > 0
  \]
\end{notebox}

% ============================================================
\section{Computing Mixed NE using Indifference}

\begin{exbox}{Computing a Mixed NE}
  Consider the following game:
  \begin{center}
  \renewcommand{\arraystretch}{1.4}
  \begin{tabular}{c|cc}
    $P_1 \setminus P_2$ & L ($y$) & R ($1-y$) \\
    \hline
    T ($x$) & \payoff{1}{-1} & \payoff{0}{2} \\
    B ($1-x$) & \payoff{0}{1} & \payoff{2}{0} \\
  \end{tabular}
  \end{center}
\end{exbox}

There is no pure strategy NE, nor is there an NE where one player plays a pure strategy and the other plays a proper mixed strategy. Thus, any NE must have $0 < x < 1$ and $0 < y < 1$.

Using the Indifference Principle for Player 1:
\begin{align*}
  U_1(T, y) &= U_1(B, y) \\
  1(y) + 0(1-y) &= 0(y) + 2(1-y) \\
  y &= 2 - 2y \implies 3y = 2 \implies y = \frac{2}{3}
\end{align*}

Using the Indifference Principle for Player 2:
\begin{align*}
  U_2(x, L) &= U_2(x, R) \\
  -1(x) + 1(1-x) &= 2(x) + 0(1-x) \\
  1 - 2x &= 2x \implies 4x = 1 \implies x = \frac{1}{4}
\end{align*}

The unique Nash Equilibrium is $\sigma_1^* = (1/4, 3/4)$ and $\sigma_2^* = (2/3, 1/3)$.

% ============================================================
\section{Algorithm for Finding All NE of a Two-Player Game}

To find all Nash Equilibria of a general two-player finite game, we can systematically search over all possible \textbf{supports}.
Let $S_1$ and $S_2$ be the finite sets of pure strategies, with $|S_1| = m_1$ and $|S_2| = m_2$.

For every possible subset $Y_1 \subseteq S_1$ and $Y_2 \subseteq S_2$ (there are $(2^{m_1}-1) \times (2^{m_2}-1)$ such pairs of supports), we check if there exists an NE $(\sigma_1^*, \sigma_2^*)$ such that $\text{supp}(\sigma_1^*) = Y_1$ and $\text{supp}(\sigma_2^*) = Y_2$.

Such a profile is an NE if and only if it satisfies the following system of linear equations and inequalities:

\begin{enumerate}
    \item \textbf{Indifference within the support (Player 1):} \\
    For some fixed $s_1^o \in Y_1$,
    \[ \sum_{s_2 \in Y_2} \sigma_2^*(s_2) u_1(s_1^o, s_2) = \sum_{s_2 \in Y_2} \sigma_2^*(s_2) u_1(s_1, s_2) \quad \forall s_1 \in Y_1 \]
    \item \textbf{Indifference within the support (Player 2):} \\
    For some fixed $s_2^o \in Y_2$,
    \[ \sum_{s_1 \in Y_1} \sigma_1^*(s_1) u_2(s_1, s_2^o) = \sum_{s_1 \in Y_1} \sigma_1^*(s_1) u_2(s_1, s_2) \quad \forall s_2 \in Y_2 \]
    \item \textbf{No profitable deviation outside the support (Player 1):}
    \[ \sum_{s_2 \in Y_2} \sigma_2^*(s_2) u_1(s_1^o, s_2) \geq \sum_{s_2 \in Y_2} \sigma_2^*(s_2) u_1(s_1, s_2) \quad \forall s_1 \notin Y_1 \]
    \item \textbf{No profitable deviation outside the support (Player 2):}
    \[ \sum_{s_1 \in Y_1} \sigma_1^*(s_1) u_2(s_1, s_2^o) \geq \sum_{s_1 \in Y_1} \sigma_1^*(s_1) u_2(s_1, s_2) \quad \forall s_2 \notin Y_2 \]
    \item \textbf{Valid probabilities (sum to 1):}
    \[ \sum_{s_1 \in Y_1} \sigma_1^*(s_1) = 1, \quad \sum_{s_2 \in Y_2} \sigma_2^*(s_2) = 1 \]
    \item \textbf{Valid probabilities (strictly positive on support):}
    \[ \sigma_1^*(s_1) > 0 \quad \forall s_1 \in Y_1, \quad \sigma_2^*(s_2) > 0 \quad \forall s_2 \in Y_2 \]
    \item \textbf{Zero probability outside support:}
    \[ \sigma_1^*(s_1) = 0 \quad \forall s_1 \notin Y_1, \quad \sigma_2^*(s_2) = 0 \quad \forall s_2 \notin Y_2 \]
\end{enumerate}

\subsection{Simplifying the Search via Dominance}

Before applying the exhaustive search, one should iteratively eliminate strictly dominated strategies to reduce the sizes of $S_1$ and $S_2$. 

\begin{exbox}{Strict Dominance by a Mixed Strategy}
  \begin{center}
  \renewcommand{\arraystretch}{1.4}
  \begin{tabular}{c|ccc}
    $P_1 \setminus P_2$ & L & C & R \\
    \hline
    T & \payoff{6}{2} & \payoff{0}{6} & \payoff{4}{4} \\
    M & \payoff{2}{12} & \payoff{4}{3} & \payoff{2}{5} \\
    B & \payoff{0}{6} & \payoff{10}{0} & \payoff{2}{2} \\
  \end{tabular}
  \end{center}
  Notice that the pure strategy $M$ for Player 1 is strictly dominated by the mixed strategy $\sigma_1 = \left(\frac{1}{2} T, 0 M, \frac{1}{2} B\right)$:
  \begin{itemize}
      \item Against L: $\frac{1}{2}(6) + \frac{1}{2}(0) = 3 > 2$
      \item Against C: $\frac{1}{2}(0) + \frac{1}{2}(10) = 5 > 4$
      \item Against R: $\frac{1}{2}(4) + \frac{1}{2}(2) = 3 > 2$
  \end{itemize}
  Since $M$ is strictly dominated, it can never be in the support of any Nash Equilibrium, and we can immediately eliminate it to simplify the game to a $2 \times 3$ game.
\end{exbox}

\end{document}
